\input{../tex-inputs/latex-leseansicht-vorspann}

\section[Rainer Maria Rilke an Arthur Schnitzler, 30. 3. 1899]{L04536 Rainer Maria Rilke an Arthur Schnitzler, 30. 3. 1899}
\nopagebreak\mylabel{L04536v}
\rehead{ }\normalsize\beginnumbering\briefempfaengerindex{Schnitzler, Arthur@\textsc{Schnitzler, Arthur}!zzzRilke, Rainer Maria@\emph{von Rainer Maria Rilke}!1899-03-301@{30. 3. 1899}|(be}
\toendnotes[C]{\smallbreak\pagebreak[2]}
\correspDesc{Versand  durch Rainer Maria Rilke am 30. 3. 1899 in Schmargendorf
\newline{}Erhalt  durch Arthur Schnitzler im Zeitraum [31. 3. 1899 – 4. 4. 1899?] in Wien}\toendnotes[C]{\smallbreak}
\Standort{CUL, Schnitzler, B 84.}
\physDesc{Brief, 1 Blatt, 4 Seiten, 2450 Zeichen
\newline{}Handschrift: schwarze Tinte, deutsche Kurrent
\newline{}Schnitzler: mit rotem Buntstift zwei Unterstreichungen }
\buchAbdrucke{\weitereDrucke{\emph{Rainer Maria Rilke und Arthur Schnitzler. Ihr Briefwechsel.}. Mit Anmerkungen herausgegeben von Heinrich Schnitzler In: \emph{Wort und Wahrheit}, Jg. 13, Nr. 4, April 1958, S. 282–298, hier S. 283–284.} }\toendnotes[C]{\smallbreak}
\pstart
           {\pb}\textsc{Schmargendorf}{ }\textsuperscript{bei}{ }\textsc{Berlin}\oindex{Schmargendorf@\textbf{Schmargendorf}|pw}\pend
           
\pstart
           \textsc{Villa Waldfrieden\oindex{Villa Waldfrieden@\textbf{Villa Waldfrieden}|pw}}\pend
           
\pstart
           am 30. März. 99.\pend
           
\pstart{}Sehr verehrter lieber Herr Doctor,\pend\vspace{0.5em}
\pstart
           wie ungeladene Gäſte, Ihrer Güte aufs Wort glaubend,{ }ſo{ }ſind meine beiden Prosabücher\pwindex{Rilke, Rainer Maria 4.\,12.\,1875 Prag – 29.\,12.\,1926 Montreux@\textsc{Rilke, Rainer Maria} (4.\,12.\,1875 Prag – 29.\,12.\,1926 Montreux)!Zwei Prager Geschichten@\strich\emph{Zwei Prager Geschichten}|pwv}\pwindex{Rilke, Rainer Maria 4.\,12.\,1875 Prag – 29.\,12.\,1926 Montreux@\textsc{Rilke, Rainer Maria} (4.\,12.\,1875 Prag – 29.\,12.\,1926 Montreux)!Am Leben hin. Novellen und Skizzen@\strich\emph{Am Leben hin. Novellen und Skizzen}|pwv} zu
               Ihnen gekommen, nicht wahr? – Oder sie kommen noch: denn ich habe den Verleger\orgindex{Adolf Bonz und Comp.@Adolf Bonz {\kaufmannsund}  Comp.|pwv} beauftragt,{ }ſie Ihnen zu{ }ſenden.
               Das kommt{ }ſo. Wenigen Menſchen gegenüber überfällt mich das Gefühl, etwas nachholen
               zu müſſen,{ }ſo{ }ſtark, wie vor Ihnen, lieber Doctor \textsc{Schnitzler}. Nach der kleinen und durch die Umſtände vielfach beſchränkten
               Zeit unſeres \label{K_L04536-1v}\edtext{Beiſammenſeins}{\lemma{\textnormal{\emph{Beisammenseins}}}\Cendnote{\textnormal{Vgl. A. S.: \emph{Tagebuch}, 16. 3. 1899.}}}\label{K_L04536-1} iſt dieſes Bedürfnis wach geworden
               und täglich iſt es gewachſen an Macht und Muth.\pend
           
\pstart
           Und nun iſt es{ }ſoweit, daſs ich meine Schande nichtmehr verbergen kann{\dots}{ }{\pb}Irgendein Vertrauen muſs ich zu Ihnen
               haben, irgendwie die kurze Stunde ausbreiten und ihre Falten breitſtreichen, und mit
               mehreren Farben den Contur dieser lieben Erinnerung füllen: Und darum geb’ ich Ihnen
               dieſe zwei Vergangenheiten, Gedanken und Dankbarkeiten aus dunkler Bücherkindheit;
               denn{ }ſie liegen \introOben{}wie\introOben{} 7 Jahre vor mir; aber ich hab’ auch
               beide lieb wie man etwas von vor{ }ſieben Jahren zu lieben anfängt mit jenem leiſeren
               Verzeihen, das eine{ }ſo{ }ſchöne Form iſt für ein{ }ſicheres Beſitzen.\pend
           
\pstart
           So will ich auch, daſs dieſe Bücher\pwindex{Rilke, Rainer Maria 4.\,12.\,1875 Prag – 29.\,12.\,1926 Montreux@\textsc{Rilke, Rainer Maria} (4.\,12.\,1875 Prag – 29.\,12.\,1926 Montreux)!Zwei Prager Geschichten@\strich\emph{Zwei Prager Geschichten}|pwv}\pwindex{Rilke, Rainer Maria 4.\,12.\,1875 Prag – 29.\,12.\,1926 Montreux@\textsc{Rilke, Rainer Maria} (4.\,12.\,1875 Prag – 29.\,12.\,1926 Montreux)!Am Leben hin. Novellen und Skizzen@\strich\emph{Am Leben hin. Novellen und Skizzen}|pwv} in Ihnen{ }ſein dürften wie Erinnerungen an flüchtige
               gemeinſame Geſpräche, die uns vor jenen jüngſten geſchehen{ }ſind. Es müſſen {\pb}nämlich{ }ſolche
                  Begegnung{[}en{]} immer adlig{ }ſein und einen Stammbaum haben: die
               erſte darf{ }ſich nicht als Beginn fühlen,{ }ſondern als Summe ferner Vorfahren. Sie{ }ſorgt dann{ }ſelber{ }ſchon, daſs das Geſchlecht nicht ausſterbe!\pend
           
\pstart
           Im Übrigen: was waren dieſe \label{K_L04536-2v}\edtext{wiener\oindex{Wien@\textbf{Wien}|pw} Tage}{\lemma{\textnormal{\emph{wiener Tage}}}\Cendnote{\textnormal{Rilke\pwindex{Rilke, Rainer Maria 4.\,12.\,1875 Prag – 29.\,12.\,1926 Montreux@\textsc{Rilke, Rainer Maria} (4.\,12.\,1875 Prag – 29.\,12.\,1926 Montreux)|pwk} war
                     zwischen 15. 3. 1899 und 18. 3. 1899
                     in Wien\oindex{Wien@\textbf{Wien}|pwk}.}}}\label{K_L04536-2} lieb und licht! Wie{ }ſchön geſteigert klangen{ }ſie in der
                  \label{K_L04536-3v}\edtext{Premiere\eventindex{Burgtheater@\textbf{Burgtheater}!Aufführung von Der Abenteurer und die Sängerin, Die Hochzeit der Sobeide, 14.5.1899@Aufführung von Der Abenteurer und die Sängerin, Die Hochzeit der Sobeide, 14.5.1899|pwv}}{\lemma{\textnormal{\emph{Premiere}}}\Cendnote{\textnormal{Aufführung von Der Abenteurer und die Sängerin, Die Hochzeit der Sobeide, 14.5.1899. In: A. S.: \emph{Kulturveranstaltungen}, \url{https://schnitzler-kultur.acdh.oeaw.ac.at/pmb207041.html}.}}}\label{K_L04536-3}\textsc{Loris\pwindex{Hofmannsthal, Hugo von 1.\,2.\,1874 Wien – 15.\,7.\,1929 Rodaun@\textsc{Hofmannsthal, Hugo von} (1.\,2.\,1874 Wien – 15.\,7.\,1929 Rodaun)|pw}}’ aus, reich und rauſchend in dieſen Verſen{ }ſchäumend über den Rand der Zeit.\pend
           
\pstart
           Freilich: mitten heraus muſste ich auf den Bahnhof. Das war gewaltſam, ein
               Aus-den-angelnheben. – Von \textsc{Prag\oindex{Prag@\textbf{Prag}|pw}} aus hab’ ich dann meine erſte Freude \label{K_L04536-4v}\edtext{an \textsc{Hofmannsthal\pwindex{Hofmannsthal, Hugo von 1.\,2.\,1874 Wien – 15.\,7.\,1929 Rodaun@\textsc{Hofmannsthal, Hugo von} (1.\,2.\,1874 Wien – 15.\,7.\,1929 Rodaun)|pw}} geſchickt}{\lemma{\textnormal{\emph{an … geschickt}}}\Cendnote{\textnormal{Der Brief lautete:
                           »Prag, Wassergasse
                        15\oindex{Vodičkova 15B@\textbf{Vodičkova 15B}|pw}.{ / }am 19. März 99.{ / }Verehrter Herr von Hofmannsthal\pwindex{Hofmannsthal, Hugo von 1.\,2.\,1874 Wien – 15.\,7.\,1929 Rodaun@\textsc{Hofmannsthal, Hugo von} (1.\,2.\,1874 Wien – 15.\,7.\,1929 Rodaun)|pw}{ / }Der Bahnnacht zum Trotz, die mich, dunkel und ungewiß, mitten aus dem Märchen\eventindex{Burgtheater@\textbf{Burgtheater}!Aufführung von Der Abenteurer und die Sängerin, Die Hochzeit der Sobeide, 14.5.1899@Aufführung von Der Abenteurer und die Sängerin, Die Hochzeit der Sobeide, 14.5.1899|pwv} von Sobeïdens Sehnsucht\pwindex{Hofmannsthal, Hugo von 1.\,2.\,1874 Wien – 15.\,7.\,1929 Rodaun@\textsc{Hofmannsthal, Hugo von} (1.\,2.\,1874 Wien – 15.\,7.\,1929 Rodaun)!Hochzeit der Sobeide@\strich\emph{Die Hochzeit der Sobeide}|pw} riß, kehre ich zu
                        Ihnen zurück von meinem Dank und meinen besten Gedanken geführt.{ / }Sie übermäßiger, unersättlicher Verschwender! Sie haben bislang Ihre Verse,
                        die rosenreifen, in die Andacht einsamer Leser und Lauscher wie in schöne
                        Schalen gestreut; gestern aber warfen Sie in göttlichem Vergeuden in das
                        Meer der Menge hinaus was der stille Stolz ihrer Gärten ist.{ / }Die Erregung des Überschütteten, Überbeschenkten, dessen Freude sich hundert
                        Hände wünscht, fiel mich aus Ihren Versen an. Die ganze Nacht trug ich einen
                        Duft in meinen Kleidern und wie ein anvertrautes Geheimnis nahm ich den
                        Nachklang von Schönheit und Schicksal mit.{ / }Sooft ich Ihre Verse gelesen habe, geschah es mir, daß ich vor einer dunkeln
                        Stelle blieb, wie man im Wald nicht weiter geht, geht irgendwo die
                        Landschaft licht oder das lichte Aug einer einsamen Wegmadonna auf. Das
                        Weitermüssen durch die Schönheit durch, zu dem die Verse von der Bühne
                        zwingen: das war der neue gewaltsame Genuß, mit dem der Abend gestern mich
                        beschenkt hat.{ / }Wenn ich Sie als den Führer oft empfand, der dunkle Worte spricht vor
                        ernsten Bildern und einen tiefern Sinn in Bäume und in Blumen senkt im
                        Weitergehen, so hab' ich Sie gestern als den Herrn gefühlt, und Ihres Wesens
                        Wille war mein Weg.{ / }Ich sage Ihnen das, lieber Herr von
                           Hofmannsthal\pwindex{Hofmannsthal, Hugo von 1.\,2.\,1874 Wien – 15.\,7.\,1929 Rodaun@\textsc{Hofmannsthal, Hugo von} (1.\,2.\,1874 Wien – 15.\,7.\,1929 Rodaun)|pw}, weil ich Ihnen an jenem kurzen Abend viel verschwieg;
                        und weil ich weiß, daß ich Ihnen seither noch näher bin; denn gestern war
                        ich der Ihre in Angst und Andacht!{ / }Ihr Rainer Maria Rilke\pwindex{Rilke, Rainer Maria 4.\,12.\,1875 Prag – 29.\,12.\,1926 Montreux@\textsc{Rilke, Rainer Maria} (4.\,12.\,1875 Prag – 29.\,12.\,1926 Montreux)|pw}{ / }Bitte, grüßen Sie Herrn Dr. Schnitzler
                        sehr!« (Hugo von
                        Hofmannsthal\pwindex{Hofmannsthal, Hugo von 1.\,2.\,1874 Wien – 15.\,7.\,1929 Rodaun@\textsc{Hofmannsthal, Hugo von} (1.\,2.\,1874 Wien – 15.\,7.\,1929 Rodaun)|pwk}, Rainer Maria Rilke\pwindex{Rilke, Rainer Maria 4.\,12.\,1875 Prag – 29.\,12.\,1926 Montreux@\textsc{Rilke, Rainer Maria} (4.\,12.\,1875 Prag – 29.\,12.\,1926 Montreux)|pwk}:
                        \emph{Briefwechsel}. Herausgegeben von Rudolf Hirsch und
                     Ingerborg Schnack. Frankfurt am Main:
                        \emph{Suhrkamp}{ }1978, S. 41–42.)}}}\label{K_L04536-4} mit ungewiſſer Adreſſe –
               kams in{ }ſeine Hand? Ich glaube immer: man kann nur mit der Freude danken. – Dann lag
               ich in \textsc{Prag\oindex{Prag@\textbf{Prag}|pw}} mit der {\pb}Influenza in einem höchſt
               fatalen Bett. Und jetzt bin ich im »Waldfrieden\oindex{Villa Waldfrieden@\textbf{Villa Waldfrieden}|pw}« wieder und habe den Wald, den weiten, wehenden, und Bilder von
                  \textsc{Monet\pwindex{Monet, Claude 14.\,11.\,1840 Paris – 5.\,12.\,1926@\textsc{Monet, Claude} (14.\,11.\,1840 Paris – 5.\,12.\,1926)|pw}} und eine Aufführung\eventindex{Theater- und Vortragssaal der Urania-Gesellschaft@\textbf{Theater- und Vortragssaal der Urania-Gesellschaft}!Aufführung von Daheim, 29.3.1899@Aufführung von Daheim, 29.3.1899|pwv} von \textsc{Maeterlincks\pwindex{Maeterlinck, Maurice 29.\,8.\,1862 Gent – 6.\,5.\,1949 Nizza@\textsc{Maeterlinck, Maurice} (29.\,8.\,1862 Gent – 6.\,5.\,1949 Nizza)|pw}} »\textsc{Intérieur\pwindex{Maeterlinck, Maurice 29.\,8.\,1862 Gent – 6.\,5.\,1949 Nizza@\textsc{Maeterlinck, Maurice} (29.\,8.\,1862 Gent – 6.\,5.\,1949 Nizza)!Intérieur@\strich\emph{Intérieur}|pw}}« in einem neuen wirklich intimen Theater\orgindex{Intimes Theater@Intimes Theater|pwv} geſehen und empfinde die Summe dieſer unaddirbaren heterogenen
               Freuden mit{ }ſeltſamer Sorgloſigkeit. Aus dieſem Gefühl heraus ko{\geminationm}t auch der Brief, der Sie grüßen will,\pend
           \pstart in herzlicher Verehrung: \spacefill\mbox{RainerMariaRilke}\pend{}\selectlanguage{ngerman}\endnumbering\briefempfaengerindex{Schnitzler, Arthur@\textsc{Schnitzler, Arthur}!zzzRilke, Rainer Maria@\emph{von Rainer Maria Rilke}!1899-03-301@{30. 3. 1899}|)be}\mylabel{L04536h}
\begin{anhang}
\end{anhang}\newcommand{\dateiname}{L04536}\newcommand{\titel}{Rainer Maria Rilke an Arthur Schnitzler, 30. 3. 1899}\newcommand{\editorInnen}{Herausgegeben von Jahnke, SelmaMüller, Martin Anton}\input{../tex-inputs/latex-leseansicht-abspann}
